% Generated draft: requires mathematical and source review.
\documentclass[11pt]{article}
\usepackage[margin=25mm]{geometry}
\usepackage{fontspec}
\setmainfont{Latin Modern Roman}
\usepackage{amsmath,amssymb,mathtools,unicode-math}
\setmathfont{Latin Modern Math}
\usepackage{enumitem,needspace,xurl,hyperref}
\setlength{\parindent}{0pt}
\setlength{\parskip}{5pt}
\setlength{\emergencystretch}{3em}
\newcommand{\N}{\mathbb{N}}
\newcommand{\Z}{\mathbb{Z}}
\newcommand{\Q}{\mathbb{Q}}
\newcommand{\R}{\mathbb{R}}
\newcommand{\C}{\mathbb{C}}
\newcommand{\F}{\mathbb{F}}
\newcommand{\A}{\mathbb{A}}
\newcommand{\PP}{\mathbb{P}}
\newcommand{\E}{\mathbb{E}}
\newcommand{\eps}{\varepsilon}
\newcommand{\dd}{\,\mathrm{d}}
\begin{document}
\section*{UnsolvedMath ID: 30000033}
\textbf{Source problem code:} catalogue entry
\subsection{Definitions and mathematical statement}
An algebraic surface is an algebraic variety of dimension two over a ground field, which the extracted question does not specify. A hypersurface in projective three-space has a presentation
\[
S=V(F)\subset\PP^3_k
\]
for a nonzero homogeneous polynomial \(F\in k[x_0,x_1,x_2,x_3]\). The source requests elementary constructions of algebraic surfaces beyond this familiar construction, chosen to illustrate a wide range of phenomena in surface theory.

This is an expository construction programme rather than a quantified existence conjecture. It does not prescribe a class of admissible constructions, a list of phenomena or invariants, or a criterion for ``elementary.'' In particular, ``beyond hypersurfaces in \(\PP^3\)'' may mean examples presented by other constructions; it need not mean surfaces that admit no hypersurface realization whatsoever. The workshop title concerns surfaces of general type with small invariants, but the extracted request does not explicitly require every example to be smooth, projective, minimal, of general type, or to satisfy numerical bounds. These hypotheses must not be supplied merely from the workshop title.

The dated assessment, checked 2026-08-21, explicitly describes the request as expository, says that the report supplies examples, and notes that there is no unique mathematical assertion against which solvedness can be tested. No particular example is reproduced in the supplied record.
\subsection{Short English statement}
Seek accessible constructions of algebraic surfaces beyond the standard hypersurface presentation in \(\PP^3\), illustrating varied behavior. This is an expository request with no specified list of phenomena or uniquely testable conjecture.
\subsection{Sources}
\begin{itemize}
\item Local UnsolvedMath registry, permanent ID 30000033, catalogue entry, \emph{Elementary Examples of Algebraic Surface Phenomena}.
\item \emph{Mini-Workshop: Classification of Surfaces of General Type with Small Invariants} (2004), Oberwolfach Reports, pp. 439--478, DOI \url{https://doi.org/10.4171/owr/2004/09}.
\item Supplied official report link: \url{https://ems.press/content/serial-article-files/45887}.
\item Supplied literature triage checked 2026-08-21, identifying the record as an expository request rather than a well-posed theorem or conjecture.
\end{itemize}
\end{document}
